\begin{lemma}
\label{editorial-source-lemma-composition-essentially-of-finite-type}
The class of ring maps which are essentially of finite type is
preserved under composition. Similarly for essentially of finite
presentation.
\end{lemma}

\begin{proof}
Let the original maps be $R\to S\to T$. Choose an $R$-algebra
$A$ of finite type and a multiplicative subset $U\subset A$,
together with the given $R$-algebra isomorphism $U^{-1}A\cong S$.
Choose a finite type $S$-algebra $B$, a multiplicative subset
$V\subset B$, and the given $S$-algebra isomorphism
$V^{-1}B\cong T$.
Write $B=S[X_1,\ldots,X_t]/H$ for an ideal $H$.
Let $H_0\subset A[X_1,\ldots,X_t]$ be its inverse image under
the coefficient map to $S[X_1,\ldots,X_t]$, and put
$C=A[X_1,\ldots,X_t]/H_0$. This is a finite type $R$-algebra.
The natural map $U^{-1}C\to B$ is an isomorphism.
To see the kernel assertion explicitly, every polynomial over
$U^{-1}A$ is $G/u$ with $G\in A[X_1,\ldots,X_t]$ and
$u\in U$, by taking a product of its finitely many coefficient
denominators. If $G/u\in H$, then $G/1=u(G/u)\in H$,
so $G\in H_0$. Conversely every polynomial in $H_0$ maps
into $H$. This proves that the localized ideal is exactly
$H$, and proves the asserted quotient isomorphism with its
original polynomial and fraction maps.

In either this construction or the finite-presentation construction
below, set
$$
W=\{c\in C:\text{the image of }c\text{ in the original }T
\text{ is a unit}\}.
$$
This is a multiplicative subset. The canonical map
$W^{-1}C\to T$ is an isomorphism as follows.
Every image of $u\in U$ in $C$ belongs to $W$, so
$C\to W^{-1}C$ extends to $B=U^{-1}C$.
For $v\in V$, write its original representative in this
localization as $d/u$ with $d\in C$ and $u\in U$.
Its image is a unit in $T$, as is the image of $u$, hence
the image of $d$ is a unit in $T$ and $d\in W$.
Thus $v$ maps to a unit of $W^{-1}C$, and there is a map
$T=V^{-1}B\to W^{-1}C$.
It sends a fraction represented as $(c/u)/(d/w)$, with
$c,d\in C$, $u,w\in U$ and $d/w\in V$, to
$cw/(ud)$. The preceding argument applied to this actual
$d/w$ gives $d\in W$, and also $u\in W$, so its
displayed denominator belongs to $W$.
Both maps are defined by the localization relations and their
composites fix $C$, the inverted images of $U$ and $V$,
and the inverted elements of $W$. They are therefore mutual
inverses. This proves the assertion for finite type.

For finite presentation choose $A$ of finite presentation over
$R$ and choose a finite presentation
$B=S[X_1,\ldots,X_t]/(F_1,\ldots,F_m)$.
The contraction $H_0$ above need not be finitely generated,
so instead perform the following finite construction, keeping
all actual coefficients and denominators.
For each $j$, write
$$
F_j=\sum_{\nu\in E_j}\frac{a_{j\nu}}{u_{j\nu}}X^\nu,
\qquad a_{j\nu}\in A,\quad u_{j\nu}\in U,
$$
where $E_j$ is the finite set of exponent tuples occurring in
the polynomial. Set
$$
d_j=\prod_{\nu\in E_j}u_{j\nu},\qquad
G_j=\sum_{\nu\in E_j}
a_{j\nu}\left(\prod_{\mu\in E_j\setminus\{\nu\}}
u_{j\mu}\right)X^\nu\in A[X_1,\ldots,X_t].
$$
The coefficient image of $G_j$ in $S[X_1,\ldots,X_t]$ is
exactly $(d_j/1)F_j$. The factor $d_j/1$ is retained and is
a unit; thus the ideals generated by the images of the $G_j$
and by the $F_j$ are equal. If $F_j=0$, use $E_j=\varnothing$,
$d_j=1$ and $G_j=0$.
Put $C=A[X_1,\ldots,X_t]/(G_1,\ldots,G_m)$.
Then $U^{-1}C\cong B$ by precisely these coefficient maps.
Moreover $C$ is finitely presented over $R$: if
$A=R[Y_1,\ldots,Y_a]/(P_1,\ldots,P_b)$, choose polynomial
lifts $\widetilde G_j\in R[Y_1,\ldots,Y_a,X_1,\ldots,X_t]$
of all the original coefficients in $G_j$.
The resulting presentation is
$$
C\cong R[Y_1,\ldots,Y_a,X_1,\ldots,X_t]
/(P_1,\ldots,P_b,\widetilde G_1,\ldots,\widetilde G_m).
$$
Indeed quotienting by the $P_i$ first gives the polynomial
ring over $A$, after which the images of the remaining listed
relations are exactly the $G_j$.
Using the same $W$ and the same inverse localization maps proves
$T\cong W^{-1}C$, as required. These arguments allow empty
variable and relation lists and multiplicative subsets containing
zero; in the latter case the corresponding localization is the
zero ring and the same universal maps apply.
\end{proof}